\begin{frame}{각 요소 하나씩 (1/2): $\mathcal{S}$, $\mathcal{A}$}
  \kb{$\mathcal{S}$, 상태 공간}
  \begin{itemize}
    \item 상태는 ``에이전트의 관점에서의 \textbf{세계의 완전한
      기술}''이어야 한다.
    \item CartPole의 상태 $[x, \dot{x}, \theta, \dot{\theta}]$ 가 4차원인
      이유: 위치 두 개만 알고 속도 두 개를 모르면 ``막대가 아직
      기울어지고 있는가, 멈춰 있는가''를 분간할 수 없다.
    \item \textbf{상태에 빠진 정보는 에이전트에게 영구히 사라진
      정보}다 --- 보상이나 전이를 통해 돌아올 통로가 없기 때문.
    \item 이 ``상태 설계''가 MDP를 세울 때 \textbf{가장 많은 시간이
      드는 일}. 연속 상태 공간도 흔하다(CartPole처럼
      $\mathcal{S} \subset \mathbb{R}^n$ 인 경우) --- 표가 아니라
      신경망으로 함수를 표현하게 된다(Week 9--11).
  \end{itemize}
  \vskip0.5em
  \kb{$\mathcal{A}$, 행동 공간}
  \begin{itemize}
    \item 이산: CartPole의 left/right, FrozenLake의 4방향.
    \item 연속: Pendulum의 $[-2, 2] \subset \mathbb{R}$ (힘의 크기
      그 자체를 고른다).
    \item 이산/연속의 차이는 곧 ``어떤 함수를 표현할가''의 문제 ---
      이산이면 상태$\times$행동 크기의 \textbf{표}, 연속이면
      \textbf{연속 출력의 신경망}(Week 9 vs Week 10--11의 분기가
      여기서 나온다).
  \end{itemize}
\end{frame}
